% Generated draft: requires mathematical and source review.
\documentclass[11pt]{article}
\usepackage[margin=25mm]{geometry}
\usepackage{fontspec}
\setmainfont{Latin Modern Roman}
\usepackage{amsmath,amssymb,mathtools,unicode-math}
\setmathfont{Latin Modern Math}
\usepackage{enumitem,needspace,xurl,hyperref}
\setlength{\parindent}{0pt}
\setlength{\parskip}{5pt}
\setlength{\emergencystretch}{3em}
\newcommand{\N}{\mathbb{N}}
\newcommand{\Z}{\mathbb{Z}}
\newcommand{\Q}{\mathbb{Q}}
\newcommand{\R}{\mathbb{R}}
\newcommand{\C}{\mathbb{C}}
\newcommand{\F}{\mathbb{F}}
\newcommand{\A}{\mathbb{A}}
\newcommand{\PP}{\mathbb{P}}
\newcommand{\E}{\mathbb{E}}
\newcommand{\eps}{\varepsilon}
\newcommand{\dd}{\,\mathrm{d}}
\begin{document}
\section*{UnsolvedMath ID: 700561}
\textbf{Source problem code:} catalogue entry
\subsection{Definitions and mathematical statement}
Let \(G\) be a finite group, and write
\[
\operatorname{Inv}(G)=\{g\in G:g^2=1,\ g\ne 1\}
\]
for its set of involutions. Suppose that \(D\subseteq\operatorname{Inv}(G)\) is a conjugacy class of \(G\) and that \(G=\langle D\rangle\). For each positive integer \(i\), define
\[
D_i=\left\{d_1\cdots d_i:\ d_1,\ldots,d_i\in D,\ d_j\ne d_k\text{ and }d_jd_k=d_kd_j\text{ whenever }j\ne k\right\}.
\]
In particular, \(D_1=D\).

Classify, up to isomorphism, the finite groups \(G\) for which there exist such a class \(D\) and a positive integer \(n\) such that \(D_1,\ldots,D_n\) are precisely the distinct conjugacy classes of involutions of \(G\). Explicitly, each \(D_i\), for \(1\le i\le n\), must be a nonempty single conjugacy class consisting of involutions, these classes must be pairwise distinct, and
\[
\operatorname{Inv}(G)=\bigsqcup_{i=1}^{n}D_i.
\]
The defining products can in general equal the identity; the stated hypothesis therefore requires \(1\notin D_i\) for \(1\le i\le n\). The source gives no explicit condition on \(D_i\) for \(i>n\), and no such additional condition is imposed here.

The supplied source states that the Fischer groups denoted there by \(F_{22}\) and \(F_{23}\) satisfy this condition with \(n=3\); this example is retained as a source assertion.
\subsection{Short English statement}
Classify the finite groups generated by one conjugacy class \(D\) of involutions for which, for some positive integer \(n\), the products of exactly \(i\) distinct pairwise commuting members of \(D\), for \(i=1,\ldots,n\), form precisely all the distinct conjugacy classes of involutions. The supplied registry labels the problem open; that label is not independently verified.
\subsection{Sources}
\begin{itemize}
\item Local UnsolvedMath registry, supplied complete source record: permanent ID \(700561\), exact source problem code \texttt{catalogue entry}, title \emph{Kourovka Notebook Problem 9.28}. The supplied status is \texttt{open}.
\item \emph{The Kourovka Notebook (Unsolved Problems in Group Theory)}, Problem 9.28, ninth issue (1984), proposed by A. A. Makhn\"ev, as attributed in the supplied record.
\item The supplied background identifies the twenty-first issue edition (2026), arXiv:1401.0300v48. The supplied collection description names E. I. Khukhro and V. D. Mazurov as editors. These bibliographic details and the stated Fischer-group examples have not been independently checked.
\end{itemize}
\end{document}
